Ein Kind sitzt auf einem Schlitten (zusammen \mSO). Es startet aus der Ruhe und gleitet einen Hang hinunter, der unten \hSO\ tiefer liegt als der Start. Wegen der Reibung kommt es unten nur mit \vSO\ an. Rechne mit \gO.

\begin{abcliste}
\abc
Wie gross ist die Lageenergie am Start (Nullniveau: unten) und die kinetische Energie unten?

\abc
Wie viel thermische Energie hat die Reibung erzeugt? Welcher Anteil der Lageenergie ist das? Wie schnell wäre der Schlitten ohne Reibung?

\abc
Wo steckt diese Energie?

\abc
Lukas sagt: „Die Reibung hat die Energie vernichtet.“ Sara sagt: „Die Reibungskraft hat die Energie verbraucht, um wirken zu können.“ Was ist an diesen Aussagen falsch?
\end{abcliste}
